\customSection[sec:dev]{Реалізація та тестування системи}{0em}

У цьому розділі описано практичну реалізацію спроєктованої системи. Розглянуто роботу основних сценаріїв — від реєстрації користувача до автоматичного збору активності із зовнішніх сервісів — та проілюстровано їх знімками реального інтерфейсу. Завершує розділ опис підходу до тестування.

\customSubSection[sec:dev_auth]{Реєстрація та автентифікація}{0em}

Вхід до системи побудовано на основі токенів. Після успішної автентифікації сервер видає короткоживучий токен доступу (JWT) та токен оновлення, який зберігається в захищеному cookie й дає змогу продовжувати сесію без повторного введення пароля. Паролі ніколи не зберігаються у відкритому вигляді — за їх хешування відповідає підсистема ASP.NET Core Identity. Щоб ускладнити автоматизовані реєстрації, форму створення облікового запису захищено механізмом перевірки Altcha: сервер відхиляє реєстрацію без коректно розв'язаної задачі. Зовнішній вигляд форми входу наведено на \textbf{рис.~\ref{fig:login}}.

\setfigure[fig:login]{0.78\textwidth}{assets/login}{Форма входу до системи}

\customSubSection[sec:dev_habits]{Керування звичками та записами}{1em}

Центральним об'єктом системи є звичка. Користувач може створювати звички, переглядати їх перелік, редагувати та видаляти. Кожна картка звички у списку показує її назву, теги, джерело автоматизації, періодичність і ціль (\textbf{рис.~\ref{fig:habits}}).

\setfigure[fig:habits]{0.9\textwidth}{assets/habits}{Список звичок користувача}

Під час створення звички користувач задає її параметри: назву й опис, тип (бінарна або вимірювана), періодичність і кількість разів за період, цільове значення з одиницею вимірювання, необов'язкову кінцеву дату та проміжний рубіж, теги, а також джерело автоматизації — GitHub, Strava або жодного (\textbf{рис.~\ref{fig:habit_create}}). Саме поле джерела автоматизації пов'язує звичку із зовнішнім сервісом: лише звички з відповідним джерелом отримуватимуть автоматичні записи.

\setfigure[fig:habit_create]{0.85\textwidth}{assets/habit-create}{Форма створення звички}

Факт виконання звички фіксується записом. Записи бувають двох видів: \textbf{ручні}, які користувач створює сам, та \textbf{автоматичні}, що породжуються системою на основі активності із зовнішніх сервісів. У переліку записів кожен елемент позначено відповідною міткою — Manual або Automation, — а для автоматичних додатково зберігається опис події, як-от тип і назва репозиторію для GitHub (\textbf{рис.~\ref{fig:entries}}).

\setfigure[fig:entries]{0.9\textwidth}{assets/entries}{Перелік записів із позначками джерела}

\customSubSection[sec:dev_integrations]{Інтеграція із зовнішніми сервісами}{1em}

Під'єднання зовнішніх сервісів винесено в розділ профілю «Integrations». Для \textbf{GitHub} користувач вводить персональний токен доступу (Personal Access Token) та, за бажанням, термін його дії (\textbf{рис.~\ref{fig:github}}). Токен зберігається в захищеному вигляді й надалі використовується системою для звернень до API GitHub від імені користувача.

\setfigure[fig:github]{0.9\textwidth}{assets/integration-github}{Налаштування інтеграції з GitHub}

Для \textbf{Strava} застосовано протокол OAuth: користувач надає системі дозвіл на доступ до своїх даних, а вона отримує пару токенів (доступу й оновлення). Оскільки токен доступу Strava має обмежений строк дії, система самостійно оновлює його за допомогою токена оновлення, коли це потрібно (\textbf{рис.~\ref{fig:strava}}).

\setfigure[fig:strava]{0.9\textwidth}{assets/integration-strava}{Налаштування інтеграції зі Strava}

\customSubSection[sec:dev_polling]{Автоматичний збір активності}{1em}

Автоматичний трекінг — ключова відмінність системи — реалізовано за допомогою фонових служб. Для кожного джерела (GitHub, Strava) працює окремий фоновий процес, який з фіксованим інтервалом (кожні дві хвилини) ініціює опитування під'єднаних облікових записів (лістинг~\ref{lst:worker}).

\begin{lstlisting}[style=nocodeframe, caption={Фоновий процес періодичного опитування GitHub}, label={lst:worker}]
protected override async Task ExecuteAsync(CancellationToken stoppingToken)
{
    using var timer = new PeriodicTimer(TimeSpan.FromMinutes(2));

    while (!stoppingToken.IsCancellationRequested)
    {
        using var scope = scopeFactory.CreateScope();
        var sender = scope.ServiceProvider
            .GetRequiredService<ISender>();
        await sender.Send(
            new PollIntegrationActivityCommand(IntegrationProvider.Github),
            stoppingToken);

        await timer.WaitForNextTickAsync(stoppingToken);
    }
}
\end{lstlisting}

Під час опитування система звертається до API сервісу, отримує активність, що сталася після моменту останньої синхронізації, і відбирає релевантні події. Для GitHub це події типу \texttt{Push} (коміти) та \texttt{PullRequest}; для кожної з них додатково завантажується текст повідомлення коміта або назва запиту на злиття, щоб сформувати змістовний опис. Для Strava подіями є завершені тренування з їх типом, назвою та дистанцією.

Відібрані активності перетворюються на події \texttt{IntegrationActivitiesDetected} і публікуються. Обробник у модулі записів зіставляє їх зі звичками користувача, відсіює вже зафіксовані за зовнішнім ідентифікатором і створює нові записи. Завдяки перевірці за зовнішнім ідентифікатором повторне опитування не призводить до дублювання — та сама подія GitHub чи Strava ніколи не буде зафіксована двічі.

\customSubSection[sec:dev_dashboard]{Візуалізація активності}{1em}

Зведену картину прогресу користувач бачить на головній панелі (\textbf{рис.~\ref{fig:dashboard}}). Вона містить теплову карту активності за місяць, лічильники кількості записів і поточної серії днів (streak), блок швидкого додавання записів та стрічку останніх подій. На наведеному знімку добре помітний результат роботи автоматичного трекінгу: у стрічці присутні записи з міткою Automation, породжені реальними подіями GitHub — комітами та запитом на злиття з репозиторію користувача.

\setfigure[fig:dashboard]{0.9\textwidth}{assets/dashboard}{Головна панель: теплова карта та стрічка активності}

\customSubSection[sec:dev_testing]{Тестування системи}{1em}

Перевірку коректності системи проведено на двох рівнях.

\textbf{Модульне тестування} спрямоване на доменний шар, де зосереджено бізнес-правила. Для модулів Habit і Tag створено окремі тестові проєкти на основі фреймворку xUnit зі збиранням покриття коду засобами Coverlet. Об'єктом перевірки є насамперед валідація об'єктів-значень та інваріанти агрегатів: наприклад, неможливість створити частоту з недодатною кількістю разів, ціль із порожньою одиницею вимірювання чи рубіж, поточне значення якого перевищує цільове. Завдяки тому, що доменний шар не залежить від інфраструктури, а помилки повертаються через шаблон результату (Result), такі перевірки виконуються без під'єднання до бази даних і зовнішніх сервісів.

\textbf{Функціональне тестування} охопило наскрізні сценарії роботи системи в реальному середовищі. Було перевірено: реєстрацію й вхід користувача, створення та редагування звичок різних типів, додавання ручних записів, під'єднання зовнішнього сервісу та — найважливіше — автоматичне формування записів. Підтвердженням коректності останнього сценарію слугує головна панель (\textbf{рис.~\ref{fig:dashboard}}): після виконання реальних комітів у під'єднаному репозиторії GitHub система самостійно, без жодних дій користувача, створила відповідні записи звички. Це демонструє, що ключова функція системи — автоматичний моніторинг активності — працює як задумано.

Для підтримання якості коду під час розробки налаштовано конвеєр безперервної інтеграції на основі GitHub Actions (файл \texttt{.github/workflows/build.yml}). Він автоматично запускається на кожне надсилання змін до репозиторію, виконує збирання рішення та прогін тестів, завдяки чому помилки виявляються на ранньому етапі, ще до злиття змін.

\customSubSection[sec:dev_testcase]{Тестовий приклад}{1em}

Для перевірки працездатності системи підготовлено набір тестових сценаріїв, що охоплюють основні функції — від реєстрації до автоматичного збору активності. Кожен сценарій містить вхідні дані, очікуваний результат та фактично отриманий результат (\tabref{tab:testcases}).

\needspace{9cm}
\begin{spacing}{1.2}
    \footnotesize
    \noindent
    \begin{xltabular}{\textwidth}{| c | >{\raggedright\arraybackslash}X | >{\raggedright\arraybackslash}X | c |}
        \caption{Тестові сценарії перевірки системи} \label{tab:testcases} \\
        \hline
        \textbf{№} & \textbf{Сценарій (вхідні дані)} & \textbf{Очікуваний результат} & \textbf{Статус} \\ \hline
        \endfirsthead

        \hline
        \textbf{№} & \textbf{Сценарій (вхідні дані)} & \textbf{Очікуваний результат} & \textbf{Статус} \\ \hline
        \endhead

        1 & Реєстрація з коректними даними & акаунт створено, видано токени, перехід на дашборд & пройдено \\ \hline
        2 & Вхід із неправильним паролем & доступ відхилено, повідомлення про помилку & пройдено \\ \hline
        3 & Створення вимірюваної звички (ціль 2, одиниця «Kms») & звичка з'являється у списку з коректними параметрами & пройдено \\ \hline
        4 & Створення звички з кількістю разів за період = 0 & відмова на рівні валідації домену & пройдено \\ \hline
        5 & Додавання ручного запису & запис із міткою Manual у переліку & пройдено \\ \hline
        6 & Коміт у під'єднаний репозиторій GitHub & протягом близько 2 хв створюється автозапис із міткою Automation & пройдено \\ \hline
        7 & Повторне опитування тієї самої події GitHub & дубль запису не створюється & пройдено \\ \hline
    \end{xltabular}
\end{spacing}

Усі заплановані сценарії завершилися успішно, що підтверджує відповідність реалізованої системи поставленим вимогам.
